\SourceLeaf{336}

\section*{[II.]}
\begin{center}VORBEREITUNGEN ZUR ALLGEMEINEN THEORIE\\DER BIQUADRATISCHEN RESTE.\end{center}
\pdfbookmark[1]{II. Vorbereitungen zur allgemeinen Theorie der biquadratischen Reste}{nachlass-biquadratic-II}
\BookRule{.15\textwidth}
\par\medskip\pdfbookmark[2]{(1.)}{nachlass-biquadratic-II-prefatory-1}\centerline{(1.)}\smallskip
Es sei $P=x+iy$, wobei weder $x$ noch $y$ eine ganze Zahl ist. Wir bezeichnen die Zahl $+1$ mit $LP,L'P,L''P,L'''P$, je nachdem $P$ im ersten, zweiten, dritten oder vierten Quadranten liegt (im ersten und zweiten Quadranten ist $[y]$ gerade, im dritten und vierten ungerade; im ersten und vierten ist $[x]$ gerade, im zweiten und dritten ungerade). In allen Fällen, in denen diese Zeichen nicht $=1$ sind, werden sie als $=0$ vorausgesetzt. Dann gelten die folgenden $24$ Relationen:

\[
\begin{array}{lll}
L(P\pm1)=L'P&L(P\pm i)=L'''P&L(P\pm1\pm i)=L''P\\
L'(P\pm1)=LP&L'(P\pm i)=L''P&L'(P\pm1\pm i)=L'''P\\
L''(P\pm1)=L'''P&L''(P\pm i)=L'P&L''(P\pm1\pm i)=LP\\
L'''(P\pm1)=L''P&L'''(P\pm i)=LP&L'''(P\pm1\pm i)=L'P
\end{array}
\]
\[
\begin{array}{lll}
LiP=L'''P&L(-P)=L''P&L(-iP)=L'P\\
L'iP=LP&L'(-P)=L'''P&L'(-iP)=L''P\\
L''iP=L'P&L''(-P)=LP&L''(-iP)=L'''P\\
L'''iP=L''P&L'''(-P)=L'P&L'''(-iP)=LP
\end{array}
\]

\par\medskip\pdfbookmark[2]{(2.)}{nachlass-biquadratic-II-prefatory-2}\centerline{(2.)}\smallskip
Mit $PP'$ oder $z$ bezeichnen wir eine Linie, die bei $P$ beginnt und bei $P'$ endet. Sie muss nicht gerade sein. Wir geben allen von $2x+2iy$ nach $2x+(2y+1)i$ gezogenen geraden Linien (wobei $x,y$ sämtliche ganzen Zahlen durchlaufen) eine positive und eine negative Seite; als positive wählen wir die rechte, als negative die linke. Mit $Tz$ bezeichnen wir die Anzahl aller Schnitte
\SourceLeaf{337}

der Linie $z$ mit den eben beschriebenen Linien, wobei diejenigen positiv gezählt werden, bei denen $z$ von der negativen auf die positive Seite übergeht, und die anderen negativ. Ferner setzen wir
\[
Tz-T(z-1)=Sz
\]

\Needspace{18\baselineskip}
($z-1$ ist eine zu $z$ parallele Linie, die vom Punkt $P-1$ nach $P'-1$ geht). Offenbar brauchen wir dem obigen Liniensystem nur noch die von $2x+1+2yi$ nach $2x+1+(2y+1)i$ gezogenen Linien hinzuzufügen und deren linke Seiten als positiv und die rechten als negativ zu betrachten, um in $Sz$ die Anzahl aller Schnitte von $z$ mit diesem doppelten System gerader Linien zu erkennen. Ferner gilt
\[
\begin{aligned}
T(-z)&=-T(z+i)\\
T(z)+T(z+i)&=[\tfrac12x]-[\tfrac12x']\\
S(z+1)&=-Sz\\
S(z+i)&=-Sz+LP+L'''P-LP'-L'''P'\\
S(z+1+i)&=Sz-LP+L'''P+LP'+L'''P'\\
Siz&=Sz-LP+LP'\\
S(-z)&=Sz-LP-L'''P+LP'+L'''P'\\
S(-iz)&=Sz+L'P-L'P'
\end{aligned}
\]

\renewcommand{\Article}[1]{\par\medskip\Needspace{4\baselineskip}\pdfbookmark[2]{Art. #1}{nachlass-biquadratic-II-art-#1}\centerline{\textbf{#1.}}\nobreak\smallskip}
\Article{1}
Wir betrachten in der Ebene zwei Arten von Punkten: erstens diejenigen, denen ganze Zahlen entsprechen; zweitens diejenigen, die durch Produkte ganzer Zahlen mit der Größe $Q=\frac{m}{2M}$ bestimmt sind. Wir können sie als Punkte erster und Punkte zweiter Ordnung unterscheiden.

\Article{2}
Verbindet man jeden Punkt zweiter Ordnung mit seinen vier Nachbarn durch gerade Linien, die wir \emph{Ligaturen} nennen werden, so wird die ganze Ebene in unendlich viele Quadrate geteilt. Die Punkte erster Ordnung liegen teils innerhalb dieser Quadrate, teils auf den Ligaturen innerhalb ihrer
\SourceLeaf{338}

Grenzen und teils auf den Grenzen der Ligaturen; Letzteres gilt, wenn sie zugleich Punkte zweiter Ordnung sind. Ist $kQ$ ein solcher Punkt, so muss, sofern $m,M$ ohne gemeinsamen Teiler und beide ungerade sind, $k$ durch $M$ teilbar sein.

\Article{3}
Bei den Ligaturen können wir auch zwischen Anfangs- und Endpunkt unterscheiden, also $PQ$ von $QP$, oder diese Unterscheidung in einigen Fällen beiseitelassen. Zwei solche Ligaturen nennen wir entgegengesetzte. Im Allgemeinen lassen sich die Ligaturen am bequemsten durch ihren Anfangs- und Endpunkt bezeichnen, die man gegebenenfalls in eine Klammer einschließen kann. Die zu einer Ligatur entgegengesetzte soll durch doppeltes Überstreichen angezeigt werden: $QP=\overline{\overline{PQ}}$.

\Needspace{13\baselineskip}
\Article{4}
Jedes der betrachteten Quadrate wird von vier solchen Ligaturen eingeschlossen
\[
\begin{gathered}
\{kQ,(k+1)Q\},\ \{(k+1)Q,(k+1+i)Q\},\\
\{(k+1+i)Q,(k+i)Q\},\ \{(k+i)Q,kQ\}\dotsfour\Omega
\end{gathered}
\]
denen es zur Rechten liegt. Dabei ist es wichtig, auf die Form der Zahl $k$ zu achten. Wir unterscheiden vier Arten von Quadraten, je nachdem $k\equiv0,1,1+i,i\ (\mathrm{mod.}\ 2)$ gilt, und verwenden die Zahlen $0,1,2,3$, die wir die jeweiligen Intensoren der Quadrate nennen.

\Article{5}
Den Ligaturen geben wir dieselben Intensoren wie den Quadraten, die ihnen zur Rechten liegen.

\Article{6}
Nun haben wir ein anderes, größeres Quadrat $\Omega'$ zu betrachten, nämlich dasjenige, das entsteht, wenn das in $4$ für $k=0$ angegebene mit $M$ multipliziert wird. Es wird also von den geraden Linien $\mu,\mu',\mu'',\mu'''$ begrenzt
\[
\{0,\tfrac12m\},\ \{\tfrac12m,\tfrac12(1+i)m\},\ \{\tfrac12(1+i)m,\tfrac12im\},\ \{\tfrac12im,0\}
\]
Es besteht aus ganzen Quadraten $\Omega$ und Teilen solcher Quadrate. Man zähle
\SourceLeaf{339}

alle Punkte erster Ordnung in seinem Innern zusammen, indem man für jeden Punkt den Intensor des Quadrats $\Omega$ nimmt, in dem er liegt. Diese Summe oder ihr kleinster Rest modulo $4$ heißt der Dezident von $M$ für den Modul $m$ und bestimmt die biquadratische Modalität von $M$ bezüglich dieses Moduls.

\Article{7}
Wir zerlegen das Quadrat $\Omega'$ auf folgende Weise in $5$ Teile. Man verbinde den Punkt $0$ mit $\frac14(1+i)(m-1)$ durch die Linie $\lambda$, die ausschließlich entlang Ligaturen innerhalb von $\Omega'$ verläuft. Es sei
\[
\tfrac12m+i\lambda=\lambda',\quad\tfrac12(1+i)m-\lambda=\lambda'',\quad\tfrac12im-i\lambda=\lambda'''
\]
Diese $4$ Linien gehen also von den Ecken des Quadrats $\Omega'$ ins Innere und enden an den vier Ecken des innersten Quadrats. Dessen Intensor ist $0$, wenn $m\equiv1\ (\mathrm{mod.}\ 2+2i)$; die Ligaturen dieses Quadrats seien $\nu,\nu',\nu'',\nu'''$.

Die $5$ Teile werden also begrenzt sein durch
\[
\begin{array}{rll}
\mathrm{I}.&\ldots&\mu,\ \lambda',\ \overline{\overline\nu},\ \overline{\overline\lambda}\\
\mathrm{II}.&\ldots&\mu',\ \lambda'',\ \overline{\overline{\nu'}},\ \overline{\overline{\lambda'}}\\
\mathrm{III}.&\ldots&\mu'',\ \lambda''',\ \overline{\overline{\nu''}},\ \overline{\overline{\lambda''}}\\
\mathrm{IV}.&\ldots&\mu''',\ \lambda,\ \overline{\overline{\nu'''}},\ \overline{\overline{\lambda'''}}\\
\mathrm{V}.&\multicolumn{2}{l}{\text{das innere Quadrat }\nu,\ \nu',\ \nu'',\ \nu'''}
\end{array}
\]
Der Dezident ist also die Aufzählung aller Punkte erster Ordnung in I. II. III. IV.

\Article{8}
Der Kürze halber soll der Intensor eines beliebigen Punktes der Intensor des Quadrats sein, in dem er liegt, und durch ein vorangestelltes $\Upsilon$ bezeichnet werden.

\Article{9}
Der Dezident ist also
\[
\Sigma\Upsilon P+\Sigma\Upsilon P'+\Sigma\Upsilon P''+\Sigma\Upsilon P'''
\]
wobei $P$ alle Punkte in I. usw. bezeichnet.

\SourceLeaf{340}

\Article{10}
Wir betrachten nun noch den Bereich $\mathrm{VI}=-i\mathrm{IV}$, der außerhalb von $\Omega'$ liegt, sich aber entlang $\mu$ an I anschließt und mit ihm zusammen den Bereich $\omega$ bildet, der aus $AA+BB$ vollständigen Quadraten besteht. Bezeichnet $\Pi$ alle ganzen und $\Pi'$ alle um $\frac12i$ vermehrten ganzen Punkte dieses Bereichs, so lässt sich leicht beweisen, dass der Dezident
\[
\begin{aligned}
={}&\Sigma\Upsilon\Pi-\Sigma\Upsilon\Pi'+\text{Anzahl aller ganzen Punkte innerhalb VI}\\
&-\text{Anzahl aller halben Punkte innerhalb VI.}
\end{aligned}
\]
